\begin{examplebox}

\textbf{\textcolor{CExample}{例1（基础）}}

\textbf{\textcolor{CExample}{题目：}} 已知三角形的三边长分别为 $x=5,\;y=6,\;z=7$，求其面积.

\textbf{\textcolor{CExample}{解题步骤：}}
\begin{description}[style=nextline, leftmargin=2.8em, labelsep=0.5em, itemsep=0.45em]
    \item[\textbf{第一步（计算平方）：}] 计算各边平方：$x^2=25,\;y^2=36,\;z^2=49$.
    \par\textbf{理由：} 需要平方值建立方程组.
    \item[\textbf{第二步（解辅助量）：}] 由 $P+Q=25,\;Q+R=36,\;R+P=49$ 解得 $P=19,\;Q=6,\;R=30$.
    \par\textbf{理由：} 解线性方程组，例如 $P = \frac{x^2+z^2-y^2}{2} = \frac{25+49-36}{2}=19$，类似可得其他.
    \item[\textbf{第三步（代入公式）：}] 计算 $PQ+QR+RP = 19\times6 + 6\times30 + 30\times19 = 114+180+570 = 864$，则 $2S = \sqrt{864} = 12\sqrt{6}$，故 $S = 6\sqrt{6}$.
    \par\textbf{理由：} 直接套用结论 $2S = \sqrt{PQ+QR+RP}$.
\end{description}

\textbf{\textcolor{CExample}{关键结论：}} \textbf{$S=6\sqrt{6}\approx 14.70$}.

\medskip
\textbf{\textcolor{CExample}{例2（稍变形）}}

\textbf{\textcolor{CExample}{题目：}} 已知三角形的三边长分别为 $\sqrt{27},\sqrt{28},\sqrt{29}$，求其面积.

\textbf{\textcolor{CExample}{解题步骤：}}
\begin{description}[style=nextline, leftmargin=2.8em, labelsep=0.5em, itemsep=0.45em]
    \item[\textbf{第一步（取平方）：}] 令 $x^2=27,\;y^2=28,\;z^2=29$.
    \par\textbf{理由：} 根式边长直接取其平方，避免开方.
    \item[\textbf{第二步（解辅助量）：}] 解 $P+Q=27,\;Q+R=28,\;R+P=29$，得 $P=14,\;Q=13,\;R=15$.
    \par\textbf{理由：} 利用公式 $P = \frac{x^2+z^2-y^2}{2} = \frac{27+29-28}{2}=14$ 等.
    \item[\textbf{第三步（代入公式）：}] 计算 $PQ+QR+RP = 14\times13 + 13\times15 + 15\times14 = 182+195+210 = 587$，则 $2S = \sqrt{587}$，$S = \frac{\sqrt{587}}{2}$.
    \par\textbf{理由：} 直接得到面积，无需计算半周长和复杂根号.
\end{description}

\textbf{\textcolor{CExample}{关键结论：}} \textbf{$S = \dfrac{\sqrt{587}}{2}\approx 12.12$}.

\medskip
\textbf{\textcolor{CExample}{例3（综合应用）}}

\textbf{\textcolor{CExample}{题目：}} 三角形的三边长分别为 $\sqrt{13},\sqrt{20},\sqrt{29}$，求其面积.

\textbf{\textcolor{CExample}{解题步骤：}}
\begin{description}[style=nextline, leftmargin=2.8em, labelsep=0.5em, itemsep=0.45em]
    \item[\textbf{第一步（取平方）：}] $x^2=13,\;y^2=20,\;z^2=29$.
    \par\textbf{理由：} 根式边长直接平方.
    \item[\textbf{第二步（解辅助量）：}] 解 $P+Q=13,\;Q+R=20,\;R+P=29$，得 $P=11,\;Q=2,\;R=18$.
    \par\textbf{理由：} $P = \frac{13+29-20}{2}=11$，$Q = \frac{13+20-29}{2}=2$，$R = \frac{20+29-13}{2}=18$.
    \item[\textbf{第三步（代入公式）：}] 计算 $PQ+QR+RP = 11\times2 + 2\times18 + 18\times11 = 22+36+198 = 256$，则 $2S = \sqrt{256}=16$，$S=8$.
    \par\textbf{理由：} 公式直接给出整数结果，验证简便.
\end{description}

\textbf{\textcolor{CExample}{关键结论：}} \textbf{$S=8$}.

\end{examplebox}
